% Concise methodology aligned with the deployed shared-encoder residual-GRU.
% Required packages: amsmath, amssymb, bm

\section{Methodology}
\label{sec:methodology}

\subsection{Problem Formulation}

We consider causal estimation of hand state and reaching intent from four
surface electromyography (sEMG) channels and four six-axis inertial measurement
units (IMUs). At time $t$, the wearable observation is
\begin{equation}
    \mathbf{x}_t=[\mathbf{e}_t,\mathbf{a}_t,\boldsymbol{\omega}_t],
\end{equation}
where $\mathbf{e}_t\in\mathbb{R}^{4}$ contains the sEMG measurements and
$\mathbf{a}_t,\boldsymbol{\omega}_t\in\mathbb{R}^{12}$ contain acceleration and
angular velocity from the four IMUs. Using only observations available at or
before $t$, the model predicts the open/close state, current Cartesian
position, continuous screen location, a dense trajectory up to 200 ms, and a
lower-rate intent trajectory up to 1 s. Long-horizon motion is represented as
a displacement from the estimated current position:
\begin{equation}
    \hat{\mathbf p}^{\mathrm{intent}}_{t+h}
    =\hat{\mathbf p}_t+\Delta\hat{\mathbf p}_{t,h}.
\end{equation}

The framework contains two stages. Stage I learns a robust gripper-state
representation using multimodal causal patch encoders. Stage II freezes and
reuses these encoders and learns current and future motion through lightweight
recurrent adapters. VIVE position and future IMU measurements are used only as
supervision and are not required during inference.

\subsection{Causal Signal Processing}

Measurements are ordered by timestamp and resampled to a uniform 100-Hz grid.
Missing values are filled using the most recent past observation for at most
20 ms; longer gaps remain invalid and are accompanied by binary validity
indicators.

Each sEMG channel is processed by a fourth-order causal Butterworth band-pass
filter with lower cutoff 20 Hz and upper cutoff
\begin{equation}
    f_{\mathrm{high}}=\min(450,0.4f_s)\ \mathrm{Hz}.
\end{equation}
Trailing root-mean-square envelopes of 20 and 50 ms produce eight sEMG
features. The short window preserves transitions, while the longer window
provides a stable activation estimate. Each of the 24 IMU channels is processed
by a trailing three-sample median filter followed by a second-order causal
15-Hz low-pass filter. The DC component is retained to preserve gravity-related
orientation information.

Channel-wise normalization statistics are calculated from valid training data
only. Standardized values are clipped to $[-20,20]$ and concatenated with their
validity indicators, producing
\begin{equation}
    \tilde{\mathbf e}_t\in\mathbb{R}^{16},\qquad
    \tilde{\mathbf i}_t\in\mathbb{R}^{48}.
\end{equation}
A frame contributes to the fused losses only when at least 75\% of both
modalities are valid. Screen coordinates are predicted in normalized form and
converted using
\begin{equation}
    \hat{\mathbf u}^{\mathrm{px}}_t
    =(W\hat u_{x,t},H\hat u_{y,t}),
\end{equation}
where $W=1920$ and $H=1080$.

\subsection{Stage I: Gripper-State Representation}

\subsubsection{Multiscale causal encoding}

sEMG and IMU are encoded separately because they describe complementary
processes. sEMG directly measures muscle recruitment, whereas IMU provides the
mechanical context needed to distinguish grasp activation from reaching,
rotation, and stabilization.

Each modality is projected to a 128-dimensional feature space. A local branch
uses causal depthwise convolutions with kernels of five and nine frames. A
context branch divides the signal into overlapping 16-frame patches with
stride four and processes them using a four-layer, four-head causal
transformer. Left padding and a causal attention mask prevent access to future
samples.

Fusion is performed independently for the local and context features. For
$r\in\{\mathrm{loc},\mathrm{ctx}\}$, the modality weights are
\begin{equation}
 \boldsymbol{\alpha}^{r}_t=
 \operatorname{softmax}\!\left(
 g_r\!\left(\operatorname{sg}
 [\mathbf h^{e,r}_t,\mathbf h^{i,r}_t]\right)\right),
\end{equation}
and the fused feature is
\begin{equation}
 \mathbf z^r_t=\phi_r\!\left(
 [\alpha^{e,r}_t\mathbf h^{e,r}_t,
  \alpha^{i,r}_t\mathbf h^{i,r}_t,
  \mathbf h^{e,r}_t-\mathbf h^{i,r}_t,
  \mathbf h^{e,r}_t\odot\mathbf h^{i,r}_t]\right).
\end{equation}
The weighted terms retain modality-specific information, while the difference
and product describe cross-modal disagreement and agreement. Stop-gradient is
applied only to the gate input; the encoders remain trainable through the
downstream losses.

\subsubsection{State prediction}

The gripper decoder combines stable context logits with high-resolution local
logits:
\begin{equation}
    \boldsymbol{\ell}^{g}_t
    =\boldsymbol{\ell}^{\mathrm{ctx}}_t
    +\sigma(\beta)\boldsymbol{\ell}^{\mathrm{loc}}_t.
\end{equation}
An additional causal sEMG-only branch reconstructs masked activation and
predicts masked state labels during Stage-I learning. This branch ensures that
the representation retains direct neuromuscular information rather than
depending entirely on IMU context. At inference, its action inputs are replaced
by learned MASK tokens.

The classifier, masked reconstruction, and kinematic supervision used to learn
this representation are defined jointly with the Stage-II losses in
Sec.~\ref{subsec:two-stage-loss}. After Stage-I checkpoint selection, the
classifier and both modality encoders are frozen.

\subsection{Stage II: State-Conditioned Motion Model}

\subsubsection{Shared features and residual adaptation}

Because the two stages may use different population statistics, Stage-II
inputs are mapped into the frozen classifier's normalization space:
\begin{equation}
 \mathbf z_c=\mathbf z_m\odot
 \frac{\boldsymbol\sigma_m}{\boldsymbol\sigma_c}
 +\frac{\boldsymbol\mu_m-\boldsymbol\mu_c}
       {\boldsymbol\sigma_c}.
\end{equation}
The frozen classifier is evaluated once, and its modality-specific context
features are reused by Stage II.

Let $\mathbf u^e_t,\mathbf u^i_t\in\mathbb R^{128}$ denote these frozen
features. Separate two-layer causal GRUs learn residual temporal adaptations:
\begin{align}
 \mathbf h^e_{1:t}&=\mathbf u^e_{1:t}
 +W^e_a\operatorname{GRU}_e(\mathbf u^e_{1:t}),\\
 \mathbf h^i_{1:t}&=\mathbf u^i_{1:t}
 +W^i_a\operatorname{GRU}_i(\mathbf u^i_{1:t}).
\end{align}
The output projections are initialized to zero, so training begins from the
frozen Stage-I representation.

The frozen open/close probability is embedded as
$\mathbf s_t=\phi_s(\operatorname{sg}(\mathbf q_t))$. IMU acts as the
mechanical base and sEMG supplies a bounded state-dependent correction:
\begin{align}
 \mathbf c_t&=\tanh\!\left(
 \phi_c([\mathbf h^e_t,\mathbf s_t])\right),\\
 \gamma_t&=\sigma\!\left(
 \phi_\gamma([\mathbf h^i_t,\mathbf h^e_t,\mathbf s_t])\right),\\
 \mathbf f_t&=\operatorname{LN}
 (\mathbf h^i_t+\gamma_t\mathbf c_t).
\end{align}
The correction layer is zero-initialized and the gate initially favors the IMU
base. Consequently, sEMG alters the motion representation only when it reduces
the supervised error.

\subsubsection{Prediction heads}

Current Cartesian position is predicted from $\mathbf f_t$ using a two-layer
MLP. The screen head predicts a probability $\pi_{t,k}$ and a bounded offset
$\boldsymbol\delta_{t,k}\in[-1/6,1/6]^2$ for each cell of an internal
$3\times3$ spatial decomposition. With fixed cell centers $\mathbf m_k$, its
continuous normalized output is
\begin{equation}
 \hat{\mathbf u}_t=\sum_{k=1}^{9}\pi_{t,k}
 \operatorname{clip}_{[0,1]}
 (\mathbf m_k+\boldsymbol\delta_{t,k}),
 \quad
 \boldsymbol\pi_t=\operatorname{softmax}
 (\boldsymbol\ell^{\mathrm{cell}}_t).
\end{equation}
This decomposition is an internal continuous-regression parameterization; it
is distinct from any display partition used only to report cell accuracy.

A shared future decoder combines $\mathbf f_t$ with the horizon encoding
\begin{equation}
 \mathbf b(\tau)=[\tau,\tau^2,\sin(\pi\tau),\cos(\pi\tau)]
\end{equation}
and predicts absolute Cartesian positions every 10 ms through 200 ms. The
long-horizon decoder additionally receives the explicit sEMG state
$\mathbf h^e_t$ and predicts displacement every 100 ms through 1 s:
\begin{equation}
 \Delta\hat{\mathbf p}_{t,h}
 =\phi_{\mathrm{intent}}
 ([\mathbf f_t,\mathbf h^e_t,\mathbf b(h)]).
\end{equation}
The explicit sEMG pathway allows early muscle activation to contribute to
intent prediction before its mechanical consequence is fully observed by the
IMU.

\subsection{Two-Stage Learning Objectives}
\label{subsec:two-stage-loss}

The stages share supervision for current and future motion, but they do not
optimize the same model. Stage I learns the multimodal representation and
gripper classifier; Stage II freezes that representation and trains only the
residual GRUs and motion heads. We therefore define the losses once and then
specify the objective used at each stage.

Let $\rho(\cdot)$ denote element-wise smooth-$L_1$, and let
$\langle\cdot\rangle$ average only valid targets. Current Cartesian position
and 6-D orientation are supervised by
\begin{equation}
 \mathcal L_{\mathrm{pos}}=\langle\rho(\hat{\mathbf p}_t-\mathbf p_t)\rangle,
 \qquad
 \mathcal L_{\mathrm{ori}}=\langle\rho(\hat{\mathbf r}^{6D}_t-
 \mathbf r^{6D}_t)\rangle.
\end{equation}
Screen error is computed after scaling the two normalized axes independently
by display width $W$ and height $H$. Stage I uses a goal-consistent screen
loss over fused, direct, and endpoint-derived predictions, whereas Stage II
uses a continuous pixel loss together with internal cell and within-cell
offset supervision:
\begin{align}
 \mathcal L_{\mathrm{screen}}^{\mathrm I}
 &=\left\langle a_t^{\mathrm I}
 [d(\hat{\mathbf u}_t,\mathbf u_t)
 +0.2d(\hat{\mathbf u}^{d}_t,\mathbf u_t)
 +0.2d(\hat{\mathbf u}^{e}_t,\mathbf u_t)]\right\rangle,\\
 \mathcal L_{\mathrm{screen}}^{\mathrm{II}}
 &=0.35\mathcal L_{\mathrm{pixel}}
 +0.15\mathcal L_{\mathrm{cell}}
 +0.10\mathcal L_{\mathrm{offset}},
\end{align}
where $d$ is pixel-scaled smooth-$L_1$,
$a_t^{\mathrm I}=0.25+2.75r_t^2$, and the Stage-II pixel term uses
$a_t^{\mathrm{II}}=0.25+3.75r_t^4$. The auxiliary cell and offset terms reduce
spatial averaging but still produce one continuous screen coordinate.

For $\mathcal H_s=\{10,20,\ldots,200\}$ ms, both stages use dense future
supervision,
\begin{equation}
 \mathcal L_{\mathrm{future}}
 =\frac{1}{|\mathcal H_s|}\sum_{h\in\mathcal H_s}
 \left\langle\rho(\hat{\mathbf p}_{t+h\mid t}-\mathbf p_{t+h})\right\rangle.
\end{equation}
Stage I additionally includes future orientation in this term. Stage II adds
arrival consistency,
\begin{equation}
 \mathcal L_{\mathrm{cons}}
 =\frac{1}{|\mathcal H_s|}\sum_{h\in\mathcal H_s}
 \left\langle\rho\!\left(
 \hat{\mathbf p}_{t+h\mid t}
 -\operatorname{sg}(\hat{\mathbf p}_{t+h\mid t+h})\right)\right\rangle,
\end{equation}
which makes a forecast agree with the detached current estimate when that time
is reached.

Stage II also predicts displacement at
$\mathcal H_l=\{100,200,\ldots,1000\}$ ms:
\begin{equation}
 \mathcal L_{\mathrm{intent}}
 =\frac{\sum_{h\in\mathcal H_l}\lambda_h
 \left\langle\rho\!\left(
 \Delta\hat{\mathbf p}_{t,h}-(\mathbf p_{t+h}-\mathbf p_t)
 \right)\right\rangle}{\sum_{h\in\mathcal H_l}\lambda_h},
 \qquad \lambda_h=e^{-h/(500\,\mathrm{ms})}.
\end{equation}
Displacements reduce sensitivity to participant-specific workspace offsets,
while the decay retains 1-s supervision without allowing the least predictable
horizons to dominate.

The state-specific Stage-I losses are class-balanced gripper cross entropy
$\mathcal L_g$, sEMG-only auxiliary cross entropy $\mathcal L_{\mathrm{aux}}$,
masked-state cross entropy $\mathcal L_{\mathrm{mask}}$, masked-sEMG mean-square
reconstruction $\mathcal L_{\mathrm{recon}}$, and within-state probability
stability $\mathcal L_{\mathrm{stable}}$. They prevent class imbalance, retain
direct neuromuscular evidence, improve dropout robustness, and suppress false
open/close switching. Stage I further applies endpoint position--orientation
loss $\mathcal L_{\mathrm{end}}$ and future IMU-difference prediction
$\mathcal L_{\mathrm{IMU}}$ as representation regularizers.

Stage II replaces those source-task auxiliaries with a training-only
sEMG-to-future-IMU loss $\mathcal L_{\mathrm{cross}}$, computed on
interval-averaged future IMU changes using the same $\lambda_h$. Its decoder
receives the sEMG representation but no IMU feature, preventing a copying
shortcut and encouraging neuromuscular features that precede motion. The
correction penalty
$\mathcal L_{\mathrm{reg}}=\langle c_{t,j}^2\rangle_{t,j}$ preserves the IMU
base when the bounded sEMG residual is not useful.

The complete objectives are
\begin{align}
 \mathcal L_{\mathrm I}={}&
 \mathcal L_g+0.20\mathcal L_{\mathrm{aux}}
 +0.15(\mathcal L_{\mathrm{mask}}+\mathcal L_{\mathrm{recon}})
 +0.03\mathcal L_{\mathrm{stable}} \nonumber\\
 &+0.20\mathcal L_{\mathrm{pos}}+0.50\mathcal L_{\mathrm{ori}}
 +0.35\mathcal L_{\mathrm{screen}}^{\mathrm I}
 +0.20\mathcal L_{\mathrm{end}} \nonumber\\
 &+0.10\mathcal L_{\mathrm{future}}^{\mathrm I}
 +0.05\mathcal L_{\mathrm{IMU}},\\[2mm]
 \mathcal L_{\mathrm{II}}={}&
 1.00\mathcal L_{\mathrm{pos}}
 +\mathcal L_{\mathrm{screen}}^{\mathrm{II}}
 +0.50\mathcal L_{\mathrm{future}} \nonumber\\
 &+0.10\mathcal L_{\mathrm{cons}}
 +0.05\mathcal L_{\mathrm{intent}}
 +0.05\mathcal L_{\mathrm{cross}}
 +0.01\mathcal L_{\mathrm{reg}}.
\end{align}
Thus, the apparent overlap is intentional: in Stage I, pose and future losses
shape transferable encoder features; in Stage II, they train the new residual
adapters and deployed motion heads. Stage II has no state, orientation, or
endpoint loss because the gripper classifier is frozen and these outputs are
not deployed. The complete network contains 3,654,603 parameters: 3,001,102
frozen Stage-I parameters and 653,501 optimized Stage-II parameters.

\subsection{Causal Inference}

The deployment pipeline applies the same causal preprocessing and retains two
seconds of history with a 300-ms warm-up. Open/close output is stabilized using
hysteresis: the state changes to closed when $P(\mathrm{close})\geq0.7$, changes
to open when $P(\mathrm{close})\leq0.3$, and is otherwise retained. The model
returns current position, screen location, 20 future positions through 200 ms,
and ten intent positions through 1 s. No VIVE measurement, future wearable
sample, masked target, or online update is used during inference.
